\ProvidesFile{Trivia.tex}[Элементарные замечания]


\subsection{Элементарные замечания}

Обычно этот раздел принято называть Мета-программирование.
Я не очень люблю этот термин, потому что как будто это не программирование, а что-то другое, за гранью возможностей языка.
Тем не менее, если вы посмотрите, то все что относят к мета программированию -- это всего лишь использование шаблонов иногда с добавлением макросов.
И по сути все эти трюки нужны для того, чтобы сделать некоторые вычисления во время компиляции, а не во время исполнения.
Потому я думаю, что название выше куда лучше будет соответствовать тому, что здесь будет обсуждаться.

Теперь, что можно обсуждать в вычислениях во время компиляции.
Как и в любом языке программирования нужно иметь достаточный набор данных и операций для исполнения.

\subsubsection{Что можно использовать}

К базовым данным можно отнести:
\begin{enumerate}
\item Литералы.
\begin{cppcode}
constexpr int x = 2; // 2 is an int literal

enum class R : int { a, b, c };
\end{cppcode}
А начиная с C++20 можно использовать литералы с плавающей точкой:
\begin{cppcode}
constexpr float x = 1.2f; // is a float literal
\end{cppcode}
Вот пример использования:
\begin{cppcode}
template<int arg>
struct A {};

template<float arg>
struct B {};

template<R arg>
struct C {};
\end{cppcode}

\item \verb"constexpr" переменные.
Такие переменные можно использовать в контексте, где значение должно быть известно во время компиляции.
Например, внутри шаблонных классов и функций:
\begin{cppcode}
constexpr int x = 1;

template<int value>
class A {};

int main() {
  const int z = 1;
  A<x> y; // x can be used here but z cannot
}
\end{cppcode}
Несмотря на то, что \verb"z" является константой, ее нельзя использовать в качестве аргумента для шаблона \verb"A", так как она не помечена \verb"constexpr".

\item Имена = адреса функций (включая инстанциированные шаблоны функций).
\begin{cppcode}
using Signature = int(int);

template<Signature function>
class A {};

int func1(int arg) {
  return 2 * arg;
}

template<class Arg>
int func2(Arg arg) {
  // do something
  return /*something*/;
}

int main () {
  A<func1> x;
  A<func2<int>> y;
  A<[](int arg) { return 2 * arg; }> z;
}
\end{cppcode}
В качестве аргумента для шаблона \verb"A" можно использовать и лямбды (начиная с c++20 ), которые не перехватывают ничего.
В целом это разумно, так как такие лямбды это просто функции без имени.

\item Начиная с C++20 можно использовать объекты структур с некоторыми ограничениями.
Я не хочу прям супер строго определять, давайте скажу более или менее то, что надо знать, а строго можно посмотреть на \href{https://en.cppreference.com/w/cpp/language/template_parameters.html#Constant_template_parameter}{cppreference}.
Вот приблизительно список ограничений:
\begin{enumerate}
\item Обязательно LiteralType.
Я не хочу говорить, что это такое, но грубо говоря это <<хорошая>> структура без \verb"union"-ов и с \verb"constexpr" конструктором и деструктором.
И это правилао рекурсивно применяется к членам и базовым классам.
Так же сюда относятся массивы от LiteralType.

\item Все поля (не статические) в классе должны быть публичные и не \verb"mutable".
Все базовые классы должны быть публичными.

\item Все поля (не статические) в классе  и все базовые классы должны удовлетворять условиям выше рекурсивно.
Так же они могут быть массивами.
\end{enumerate}
Например можно писать так
\begin{cppcode}
struct S {
  int x = 0;
};

template<S arg>
struct C {};

int main() {
  C<S{}> x;
  C<S{1}> y;
}
\end{cppcode}

\item Имена классов и структур.
\begin{cppcode}
class A {};
class B {};

template<class... Args>;
class C {};

using List = C<A, B, int, double>;
\end{cppcode}

\item Имена шаблонов и их параметры.
\begin{cppcode}
template<class List>
struct SizeOfImpl {
  /*some error message*/
}

template<template<class...> class List, class... Args>
struct SizeOfImpl<List<Args...>> {
  static constexpr int value = sizeof...(Args);
}

template<class List>
constexpr int SizeOf = SizeOfImpl<List>::value;
\end{cppcode}
На 6-7 строчках используется шаблонный класс \verb"List" в качестве аргумента для структуры.
Это позволяет использовать не просто типы, а типы параметризованные значениями, что по сути является функцией на уровне типов.
\end{enumerate}

\subsubsection{Что нельзя использовать}

\begin{enumerate}
\item Литералы с плавающей точкой например, \verb"1.2".
Если вы не используете все еще c++20.

\item Строковые литералы, например, \verb$"abc"$.

\item Пространства имен -- namespace-ы.
\end{enumerate}

\paragraph{До C++20}

В мире до C++20 в языке невозможно написать что-то такое
\begin{cppcode}
template<double arg>
struct A {};

template<const char* str>
struct B {};

int main() {
  A<1.2> x;
  B<"123"> y;
}
\end{cppcode}
В обоих случаях вы получите сообщение об ошибке.
С другой стороны, если вы используете C++20, то литералы с плавающей точкой становятся доступны по умолчанию.
А строковые литералы можно заставить работать ровно так как написано выше с помощью небольшой дополнительной работы.

В мире до C++20 эту проблему можно вылечить костылями, а именно: завернуть константу в лямбду так:
\begin{cppcode}
using Sig1 = double();

template<Sig1 function>
struct A {};

using Sig2 = const char*();

template<Sig2 function>
struct B {};

int main() {
  A<[](){ return 1.2; }> x;
  B<[](){ return "123"; }> y;
}
\end{cppcode}
Выглядит это ужасно и для улучшения читаемости это безобразие можно завернуть в макрос (это тот случай, когда без макросов обойтись нельзя, только в таких случаях надо пользоваться макросами):
\begin{cppcode}
#define CONST(arg) [](){ return arg; }

template<class T>
using TypeOf = T();

template<TypeOf<double> function>
struct A {};

template<TypeOf<const char*> function>
struct B {};

int main() {
  A<CONST(1.2)> x;
  B<CONST("123")> y;
}
\end{cppcode}
Читаемость сомнительная, но лучше я не могу предложить сейчас.
Так же важная проблема -- типы, которые выглядят одинаково, будут на самом деле разными.
Например код
\begin{cppcode}
std::cout << "are same = "<< std::is_same_v< A<CONST(1.2)>, A<CONST(1.2)>> << 'n';
\end{cppcode}
выдаст
\begin{cppcode}
are same = 0
\end{cppcode}
Потому что каждая лямбда имеет свой уникальный тип.
Для литералов с плавающей точкой есть еще трюк с \verb"std::bit_cast", который доступен с C++20.
Но в C++20 есть просто на много более удобный подход, который будет рассмотрен ниже.

\paragraph{Начиная с C++17}

Из коробки в C++17 все еще нельзя написать так
\begin{cppcode}
template<const char* arg>
struct A {};

int main() {
  A<"123"> x;
}
\end{cppcode}
Но можно сделать промежуточный объект \verb"StringLiteral", который поможет сохранить синтаксис использования как мы ожидаем.
\begin{cppcode}
template<size_t N>
struct StringLiteral {
  char data[N]{};
  
  constexpr StringLiteral(const char (&str)[N]) {
    std::copy_n(str, N, data);
  }
};

template<StringLiteral arg>
struct A {};

int main() {
  A<"123"> x;
}
\end{cppcode}
И теперь код вида
\begin{cppcode}
std::cout << "are same = "<< std::is_same_v<A<"123">, A<"123">> << '\n';
\end{cppcode}выдаст
\begin{cppcode}
are same = 1
\end{cppcode}
И не надо никаких танцев с бубнами.
В целом это небольшая цена за полный функционал.

Обратите внимание на строчки
\begin{cppcode}
template<StringLiteral arg>
struct A {};
\end{cppcode}
Здесь не указан шаблонный параметр для \verb"StringLiteral".
Это связано с тем, что начиная с C++17 есть Class Template Argument Deduction (CTAD).
Вы можете использовать имя шаблона, не указывая его параметры, если эти параметры можно вычленить из инициализирующего выражения.
В примере выше
\begin{cppcode}
A<"123"> x;
\end{cppcode}
Мы инициализируем параметр строкой \verb$"123"$.
Технически это массив длины $4$ (включая терминальный символ \verb"'\0'").
Значит в данном случае инстанциируется класс
\begin{cppcode}
struct A<StringLiteral<4>> {};
\end{cppcode}

\subsubsection{Операции}

Если говорить про операции, то тут можно совершать операции над типами, а можно над данными (литералами и константами).

\paragraph{Операции над данными}

Во время компиляции у вас есть следующие способы реализовать такие операции:
\begin{enumerate}
\item С помощью шаблонов.
Классический подход заключается в двух этапном применении шаблона:
\begin{cppcode}
constexpr int x = 1;

template<int arg>
struct TwiceImpl {
  static constexpr int value = 2 * arg;
};

template<int arg>
constexpr int Twice = TwiceImpl<arg>::value;

constexpr int twice = Twice<x>;
\end{cppcode}
Можно написать еще проще:
\begin{cppcode}
constexpr int x = 1;

template<int arg>
constexpr int Twice = 2 * arg;

constexpr int twice = Twice<x>;
\end{cppcode}
Однако, первый подход лучше, если вы хотите передать \verb"Twice" как аргумент в шаблон.
То есть вы хотите передать метод через шаблонный аргумент.
В этом случае нельзя предать шаблонную переменную в качестве аргумента.
То есть такая попытка
\begin{cppcode}
template<int arg>
constexpr int Twice = 2 * arg;

template<template<int> int func> // compilation error
struct A {};
\end{cppcode}
приведет к ошибке компиляции.
В этом случае единственным возможным вариантом будет использовать промежуточный класс:
\begin{cppcode}
constexpr int x = 1;

template<int arg>
struct Twice {
  static constexpr int value = 2 * arg;
};

template<int arg>
constexpr int TwiceValue = Twice<arg>::value;

template<template<int> class func>
struct A {};

int main() {
  A<Twice> x;
}
\end{cppcode}
В этом случае изначальную структуру лучше назвать без суффикса \verb"Impl", а переменную назвать значением этой <<функции>>.
В STL принято структуру называть так как положено, а в имени переменной добавлять суффикс \verb"_v" к имени структуры.

\item С помощью \verb"constexpr" функций.
Почти любую функцию можно исполнять в константном контексте.
Для этого надо лишь пометить ее как \verb"constexpr" и если компилятор ругнется на что-то, то вы узнаете, почему ее нельзя было использовать вместе с \verb"constexpr".
\begin{cppcode}
constexpr int twice(int arg) {
  return 2 * arg;
}

template<int arg>
class A {};

int main() {
  A<twice(3)> x;
}
\end{cppcode}

\item С помощью макросов.
НИКОГДА так НЕ делайте.
\begin{cppcode}
#define twice(x) 2 * x

template<int arg>
class A {};

int main() {
  A<twice(2)> x;
}
\end{cppcode}
Вот так НИКОГДА НЕ делайте.
\end{enumerate}

\paragraph{Операции над типам}

Во время компиляции у вас есть следующие способы реализовать такие операции:
\begin{enumerate}
\item С помощью шаблонов.
Классический двухступенчатый подход:
\begin{cppcode}
template<class...>
class List;

template<class List>
struct VectorizeImpl {
  /*react to errors*/
};

template<template<class...> class List, class... Elements>
struct VectorizeImpl <List<Elements...>> {
  using type = List<std::vector<Elements>...>;
};

template<class List>
using Vectorize = typename VectorizeImpl<List>::type;

using MyList = List<int, double>;

using VecList = Vectorize<MyList>;
//VecList = List<std::vector<int>>, std::Vector<double>>;
\end{cppcode}
Причина такого подхода заключается в том, что \verb"using" нельзя специализировать.
А потому приходится вводить промежуточную структуру.
Подобный трюк используется когда надо сделать частичную специализвацию для функции.

\item С помощью макросов.
НИКОГДА так НЕ делайте.
\begin{cppcode}
#define double char

template<double arg>
class A {};

int main() {
  A<'2'> x;
}
\end{cppcode}
Я лично найду вас и казню, если вы так сделаете.
\end{enumerate}

Самое полезное для нас -- это создание списков из классов, литералов и функций, и умение выполнять над ними различную работу.
Я хочу в этом разделе обсудить в подход к работе над списками классов и элементов (фиксированного и разного типов).
Такие списки полезны при имплементации стирающих типов и в целом при выполнении содержательной работы во время компиляции.
А позже обсудить различные механизмы, использующие такие списки.
Например, \verb"compile for loop" или \verb"for each" версия.

\subsubsection{Вывод информации}
\label{section::ErrorPrint}

\paragraph{Простой вывод}

Если вам нужно получать информацию о данных во время компиляции, то единственный способ -- это получать данные в виде сообщений об ошибках.
Самый простой способ получить информацию о типах или переменных -- это использование неопределенного шаблона.
\begin{cppcode}
template<class>
struct TD;

template<auto>
struct VD;

int main() {
  constexpr int x = 5;
  TD<decltype(x)> a;
  VD<x> b;
}
\end{cppcode}
Вывод об ошибках будет (зависит от компилятора)
\begin{cppcode}
implicit instantiation of undefined template 'TD<const int>'
implicit instantiation of undefined template 'VD<5>'
\end{cppcode}
Названия \verb"TD" и \verb"VD" это сокращения от \verb"Type Detector" и \verb"Value Detector".
Первый заметный минус этого подхода -- надо писать тип и еще имя ненужной переменной.
Можно сделать синтаксис более приятным с помощью шаблонных переменных.
Можно попробовать так:
\begin{cppcode}
template<class>
struct TD;

template<auto>
struct VD;

template<class T>
constexpr TD<T> PrintT;

template<auto T>
constexpr VD<T> PrintV;

int main() {
  constexpr int x = 5;
  PrintT<decltype(x)>;
  PrintV<x>;
}
\end{cppcode}
Тогда сообщения об ошибках будут такого сорта:
\begin{cppcode}
constexpr variable cannot have non-literal type 'const TD<const int>'
implicit instantiation of undefined template 'TD<const int>'

constexpr variable cannot have non-literal type 'const VD<5>'
implicit instantiation of undefined template 'VD<5>'
\end{cppcode}

\paragraph{MSVC зависимая часть}

Мы видим лишнее сообщение об ошибке, что нельзя определять \verb"constexpr" переменную не литерального типа.
Оказывается, чтобы убрать эту ошибку достаточно сделать следующий трюк: нужно ввести промежуточный тип, который не даст развернуться анализу шаблона дальше.
Честно говоря, я не на столько знаю тонкости обработки шаблонного кода языком, чтобы сказать, какие правила в точности тут работают.
\begin{cppcode}
template<class>
struct TD;

template<auto>
struct VD;

template<class T>
struct Type : TD<T> {};

template<auto T>
struct Value : VD<T> {};

template<class T>
constexpr Type<T> PrintT;

template<auto T>
constexpr Value<T> PrintV;

int main() {
  constexpr int x = 5;
  PrintT<decltype(x)>;
  PrintV<x>;
}
\end{cppcode}
Теперь сообщения об ошибках лучше
\begin{cppcode}
implicit instantiation of undefined template 'TD<const int>'

implicit instantiation of undefined template 'VD<5>'
\end{cppcode}
Так можно проверять тип и значение любого сколь угодно сложного выражения.
